\scanpage{6}

\begin{center}
U N I V E R S I T Y\\
OF SCIENCE AND TECHNOLOGY\\
OF LANGUEDOC\\
\textasciicircum\textasciicircum\textasciicircum\textasciicircum\textasciicircum\textasciicircum\textasciicircum\textasciicircum\textasciicircum\textasciicircum\textasciicircum\textasciicircum\textasciicircum
\end{center}

Montpellier, 14 October 1983

Student Records Office

Distribution : List A\\
MA/gs \hfill Number 4375

\begin{center}
NOTICE OF DEFENCE OF A THIRD-CYCLE DOCTORAL THESIS

\uuline{GENERAL AND APPLIED ECOLOGY AQUATIC ECOLOGY OPTION}
\end{center}

\uline{SURNAME and First name of the candidate} : MONTI Dominique

\uline{Title of the thesis} : ``Research on the optimisation of the sampling design used for mapping benthic communities. Application to the benthic macrofauna of the Pierre Blanche lagoon (Palavas, Hérault, France).

\uline{Scheduled date,-time and place of defence} : On 2 November 1983 at 15:00,\\
Lecture room number 9 - University of Science and Technology of Languedoc

\uline{Name of the Thesis supervisor} : Mister AMANIEU

\uline{Designation of the Laboratory where the thesis is being prepared} : Marine Hydrobiology Laboratory - University of Science and Technology of Languedoc.

\uline{Analysis of the work} : This reflection concerns the problem of the choice of the number and spatial organisation of sampling points in ecology. What effort must one expend for a fixed level of precision ? And where must one sample ?

This work is based on carrying out different subsamplings of initial information assumed to be maximal. The different subsamplings vary in the number and spatial organisation of the sampling points ; and they are all carried out from the same list of objects.

Each subsampling used will give rise to a map which will subsequently be compared with the so-called ``reference'' map ; it will remain to assess in the light of these results whether the maps obtained differ from the reference, in what proportions and in what manner.

The theoretical results are applied to the study of the benthic macrofauna of a coastal lagoon (the Pierre Blanche lagoon at Palavas, France).

\uline{Approximate duration of the research that led to the thesis} : 3 years

\uline{Composition of the Proposed examining committee} : \begin{tabular}[t]{ll}
Misters & AMANIEU M.\\
& BLANC F.\\
& FRONTIER S.\\
& GRANDJOUAN G.\\
& LASSERRE G.
\end{tabular}

\begin{flushright}
For the Secretary General,\\
The Principal Officer for\\
School and University Administration,

\unsure{A. Gentil}

A. GENTIL
\end{flushright}

\margin{5}
